% Part: computability
% Chapter: computability-theory
% Section: rice-theorem

\documentclass[../../../include/open-logic-section]{subfiles}

\begin{document}

\olfileid{cmp}{thy}{rce}
\olsection{مبرهنة رايس}

إذا رجعنا إلى البرهان وجدنا أن ما تختص به $\fn{Tot}$ لا يدخل
في حجة \olref[tot]{prop:total}. وإنما جعلنا ${\OLId{latin-ordinary}{h}}({\OLId{latin-ordinary}{x}},{\OLId{latin-ordinary}{y}})$ موافقة
للدالة الثابتة ${\OLId{latin-ordinary}{j}}({\OLId{latin-ordinary}{y}})={\OLMathNumeral{0}}$ متى كان ${\OLId{latin-ordinary}{x}}$ في ${\OLId{latin-looped}{K}}$، لا غير؛ وكان
يجوز أن نجعلها موافقة لأي دالة جزئية قابلة
للحساب في تلك الحالة. ومن هذه الملاحظة تنشأ مبرهنة أعم، حاصلها على وجه
التقريب أن كل خاصية غير تافهة للدوال القابلة للحساب
يمتنع قرارها.

ونستحضر أن $\cfind{{\OLMathNumeral{0}}}$، و$\cfind{{\OLMathNumeral{1}}}$، و$\cfind{{\OLMathNumeral{2}}}$،~\dots هو تعدادنا
المعياري للدوال الجزئية القابلة للحساب.

\begin{thm}[مبرهنة رايس]
  لتكن ${\OLId{latin-looped}{C}}$ أي مجموعة من الدوال الجزئية القابلة للحساب، ولتكن ${\OLId{latin-looped}{A}} =
  \Setabs{{\OLId{latin-ordinary}{n}}}{\cfind{{\OLId{latin-ordinary}{n}}} \in {\OLId{latin-looped}{C}}}$. إذا كانت ${\OLId{latin-looped}{A}}$ قابلة للحساب، فإما أن
  تكون ${\OLId{latin-looped}{C}}$ خالية، وإما أن تكون ${\OLId{latin-looped}{C}}$ مجموعة جميع الدوال الجزئية
  القابلة للحساب.
\end{thm}

نسمّي ${\OLId{latin-looped}{A}}$ \emph{مجموعة فهارس} إذا لم تميز بين فهرسين لدالة واحدة:
فمتى كان ${\OLId{latin-ordinary}{n}}$ و${\OLId{latin-ordinary}{m}}$ «يحسبان» الدالة نفسها، انتميا جميعًا
إلى ${\OLId{latin-looped}{A}}$ أو خرجا منها جميعًا. وظاهر أن
المجموعة~${\OLId{latin-looped}{A}}$ الواردة في المبرهنة على هذه الصفة. وبالعكس، متى كانت ${\OLId{latin-looped}{A}}$
مجموعة فهارس، وجعلنا ${\OLId{latin-looped}{C}}$ مجموعة الدوال التي تحسبها، كان
${\OLId{latin-looped}{A}} = \Setabs{{\OLId{latin-ordinary}{n}}}{\cfind{{\OLId{latin-ordinary}{n}}} \in {\OLId{latin-looped}{C}}}$.

\begin{explain}
فتكون عبارة مبرهنة رايس بهذا الاصطلاح: لا مجموعة فهارس غير تافهة
تقبل القرار. ولا يستبين المراد حتى نميز
\emph{البرامج}، أيا كانت لغة كتابتها، من الدوال
التي تحسبها. فللبرامج، من حيث هي تراكيب رمزية،
أسئلة يقبل جوابها الحساب: أفي البرنامج أكثر من \OLClassicalDigits{150} رمزًا؟ أيتجاوز \OLClassicalDigits{22}
سطرًا؟ أتوجد فيه عبارة «while»؟ أترد «hello world» وسيطةً
لعبارة «print»؟ وأما مبرهنة رايس فتمنع القرار الحسابي لكل سؤال غير تافه عن
\emph{سلوك} البرنامج. ومثاله: أيتوقف
البرنامج عند الدخل ${\OLMathNumeral{0}}$؟ أيتوقف عند دخل ما؟ أيعطي في حال من الأحوال عددًا زوجيًا؟
\end{explain}

\begin{proof}[برهان مبرهنة رايس]
لنفرض أن ${\OLId{latin-looped}{C}}$ غير خالية، وأنها لا تشمل كل الدوال الجزئية القابلة
للحساب، ولنجعل ${\OLId{latin-looped}{A}}$ مجموعة فهارس دوال~${\OLId{latin-looped}{C}}$. ونثبت أنه
لو أمكن حساب ${\OLId{latin-looped}{A}}$ لأمكن حل مسألة التوقف؛ فيلزم أن ${\OLId{latin-looped}{A}}$ لا
تقبل الحساب.

ولا يضيق عموم البرهان إذا فرضنا الدالة ${\OLId{latin-ordinary}{f}}$ التي لا تُعرّف
في أي موضع خارجة عن ${\OLId{latin-looped}{C}}$؛ وإلا أقمنا متممة ${\OLId{latin-looped}{C}}$ مقامها في الحجة
الآتية. ونأخذ ${\OLId{latin-ordinary}{g}}$ من~${\OLId{latin-looped}{C}}$. فلو استطعنا قرار
${\OLId{latin-looped}{A}}$ لميزنا الفهارس الحاسبة لـ~${\OLId{latin-ordinary}{f}}$ من الحاسبة لـ~${\OLId{latin-ordinary}{g}}$،
واستعملنا هذا التمييز في حل مسألة التوقف.

وبيان ذلك أن الدوال الجزئية الشاملة القابلة للحساب تتيح لنا
تعريف الدالة
\[
{\OLId{latin-ordinary}{h}}({\OLId{latin-ordinary}{x}},{\OLId{latin-ordinary}{y}}) \simeq
\begin{cases}
\text{غير معرّفة} & \text{إذا كانت $\cfind{{\OLId{latin-ordinary}{x}}}({\OLId{latin-ordinary}{x}}) \fundefined$} \\
{\OLId{latin-ordinary}{g}}({\OLId{latin-ordinary}{y}}) & \text{خلاف ذلك.}
\end{cases}
\]
فإذا أردنا حساب ${\OLId{latin-ordinary}{h}}$ بدأنا بحساب $\cfind{{\OLId{latin-ordinary}{x}}}({\OLId{latin-ordinary}{x}})$؛ فإن انتهى
شرعنا في حساب~${\OLId{latin-ordinary}{g}}({\OLId{latin-ordinary}{y}})$، فإن انتهت \emph{هذه} العملية كذلك أخرجنا
قيمتها. والصياغة الأدق أن نكتب
\[
{\OLId{latin-ordinary}{h}}({\OLId{latin-ordinary}{x}},{\OLId{latin-ordinary}{y}}) \simeq \Proj{{\OLMathNumeral{2}}}{{\OLMathNumeral{0}}}({\OLId{latin-ordinary}{g}}({\OLId{latin-ordinary}{y}}),\fn{Un}({\OLId{latin-ordinary}{x}},{\OLId{latin-ordinary}{x}})).
\]
حيث $\Proj{{\OLMathNumeral{2}}}{{\OLMathNumeral{0}}}({\OLId{latin-ordinary}{z}}_{\OLMathNumeral{0}},{\OLId{latin-ordinary}{z}}_{\OLMathNumeral{1}})={\OLId{latin-ordinary}{z}}_{\OLMathNumeral{0}}$ دالة الإسقاط الثنائية التي تعيد
الوسيطة رقم~${\OLMathNumeral{0}}$، وهي قابلة للحساب.

فـ${\OLId{latin-ordinary}{h}}$ تركيب من دوال جزئية قابلة للحساب؛ ولا يتعين الطرف الأيمن
بالقيمة~${\OLId{latin-ordinary}{g}}({\OLId{latin-ordinary}{y}})$ إلا إذا تعينت $\fn{Un}({\OLId{latin-ordinary}{x}},{\OLId{latin-ordinary}{x}})$ و${\OLId{latin-ordinary}{g}}({\OLId{latin-ordinary}{y}})$
جميعًا، ويتعين كلما تعينتا.

فلنثبت~${\OLId{latin-ordinary}{x}}$: إن كانت $\cfind{{\OLId{latin-ordinary}{x}}}({\OLId{latin-ordinary}{x}})$ غير معرّفة كانت
${\OLId{latin-ordinary}{h}}({\OLId{latin-ordinary}{x}},{\OLId{latin-ordinary}{y}})$ غير معرّفة لكل~${\OLId{latin-ordinary}{y}}$؛ وإذا كانت $\cfind{{\OLId{latin-ordinary}{x}}}({\OLId{latin-ordinary}{x}})$ معرّفة كان
${\OLId{latin-ordinary}{h}}({\OLId{latin-ordinary}{x}},{\OLId{latin-ordinary}{y}}) \simeq {\OLId{latin-ordinary}{g}}({\OLId{latin-ordinary}{y}})$. فلكل~${\OLId{latin-ordinary}{x}}$ ثابتة، لا تخلو
${\OLId{latin-ordinary}{h}}({\OLId{latin-ordinary}{x}},{\OLId{latin-ordinary}{y}})$ من أن توافق ${\OLId{latin-ordinary}{f}}$ أو توافق ${\OLId{latin-ordinary}{g}}$؛ فقرار ما إذا كانت
$\cfind{{\OLId{latin-ordinary}{x}}}({\OLId{latin-ordinary}{x}})$ معرّفة هو بعينه تقرير أي الحالتين حاصل. لكن هذا يكافئ
تقرير ما إذا كانت ${\OLId{latin-ordinary}{h}}_{\OLId{latin-ordinary}{x}}({\OLId{latin-ordinary}{y}}) \simeq {\OLId{latin-ordinary}{h}}({\OLId{latin-ordinary}{x}},{\OLId{latin-ordinary}{y}})$ تنتمي إلى~${\OLId{latin-looped}{C}}$، ولو كانت ${\OLId{latin-looped}{A}}$
قابلة للحساب لأمكننا فعل ذلك تمامًا.

وتحرير ذلك أن ${\OLId{latin-ordinary}{h}}$، لكونها جزئية قابلة للحساب، هي الدالة
$\cfind{{\OLId{latin-ordinary}{e}}}[{\OLMathNumeral{2}}]$ عند فهرس ما~${\OLId{latin-ordinary}{e}}$. وتفيد مبرهنة ${\OLId{latin-ordinary}{s}}$-${\OLId{latin-ordinary}{m}}$-${\OLId{latin-ordinary}{n}}$ بوجود دالة
عودية بدائية~${\OLId{latin-ordinary}{s}}$ بحيث، لكل~${\OLId{latin-ordinary}{x}}$،
$\cfind{{\OLId{latin-ordinary}{s}}({\OLId{latin-ordinary}{e}},{\OLId{latin-ordinary}{x}})}({\OLId{latin-ordinary}{y}})\simeq {\OLId{latin-ordinary}{h}}_{\OLId{latin-ordinary}{x}}({\OLId{latin-ordinary}{y}})$. والآن، لكل ${\OLId{latin-ordinary}{x}}$، إذا كانت
$\cfind{{\OLId{latin-ordinary}{x}}}({\OLId{latin-ordinary}{x}}) \fdefined$ فإن $\cfind{{\OLId{latin-ordinary}{s}}({\OLId{latin-ordinary}{e}},{\OLId{latin-ordinary}{x}})}$ هي الدالة نفسها~${\OLId{latin-ordinary}{g}}$،
فيكون ${\OLId{latin-ordinary}{s}}({\OLId{latin-ordinary}{e}},{\OLId{latin-ordinary}{x}})$ في ${\OLId{latin-looped}{A}}$. وأما إذا كانت $\cfind{{\OLId{latin-ordinary}{x}}}({\OLId{latin-ordinary}{x}})
\uparrow$ فإن $\cfind{{\OLId{latin-ordinary}{s}}({\OLId{latin-ordinary}{e}},{\OLId{latin-ordinary}{x}})}$ هي الدالة نفسها~${\OLId{latin-ordinary}{f}}$، فلا يكون
${\OLId{latin-ordinary}{s}}({\OLId{latin-ordinary}{e}},{\OLId{latin-ordinary}{x}})$ في~${\OLId{latin-looped}{A}}$. وبعبارة أخرى، لكل~${\OLId{latin-ordinary}{x}}$، يكون ${\OLId{latin-ordinary}{x}} \in {\OLId{latin-looped}{K}}$ إذا وفقط إذا
كان ${\OLId{latin-ordinary}{s}}({\OLId{latin-ordinary}{e}},{\OLId{latin-ordinary}{x}}) \in {\OLId{latin-looped}{A}}$. فلو قبلت ${\OLId{latin-looped}{A}}$ الحساب لقبلته ${\OLId{latin-looped}{K}}$، وهو
محال. فثبت عدم قابلية ${\OLId{latin-looped}{A}}$ للحساب.
\end{proof}

ولهذه المبرهنة تطبيقات واسعة، منها النتيجة المباشرة الآتية.
\begin{cor}
المجموعات الآتية غير قابلة للقرار.
\begin{enumerate}
\item $\Setabs{{\OLId{latin-ordinary}{x}}}{\text{العدد ${\OLMathNumeral{17}}$ في مدى $\cfind{{\OLId{latin-ordinary}{x}}}$}}$
\item $\Setabs{{\OLId{latin-ordinary}{x}}}{\text{$\cfind{{\OLId{latin-ordinary}{x}}}$ ثابتة}}$
\item $\Setabs{{\OLId{latin-ordinary}{x}}}{\text{$\cfind{{\OLId{latin-ordinary}{x}}}$ كلية}}$
\item $\Setabs{{\OLId{latin-ordinary}{x}}}{\text{متى كان ${\OLId{latin-ordinary}{y}} < {\OLId{latin-ordinary}{y}}'$ وكانت $\cfind{{\OLId{latin-ordinary}{x}}}({\OLId{latin-ordinary}{y}}) \fdefined$،
    فإذا كانت $\cfind{{\OLId{latin-ordinary}{x}}}({\OLId{latin-ordinary}{y}}') \fdefined$ كان $\cfind{{\OLId{latin-ordinary}{x}}}({\OLId{latin-ordinary}{y}}) < \cfind{{\OLId{latin-ordinary}{x}}}({\OLId{latin-ordinary}{y}}')$}}$
\end{enumerate}
\end{cor}

\begin{proof}
  إذ كل واحدة من هذه مجموعة فهارس غير تافهة.
\end{proof}

\end{document}
